\documentclass[12pt]{article}
\usepackage{notes-project}

\begin{document}

\begin{center}
  {\LARGE\bfseries Lecture 1: Incompleteness of Rational Numbers}\\[0.25em]
  {\large Mathematical Analysis}\\[0.15em]
\end{center}

\section{Numbers}

\begin{dfn}[Natural Numbers]{natural-number}
Let \(S\) denote the \emph{successor}. The \emph{natural numbers} are
inductively generated by
\[
  \NN \mathrel{:=} 0 \mid S(n:\NN).
\]
These constructors exhaust \(\NN\).
\end{dfn}

\begin{thm}[Natural Number Induction]{natural-number-induction}
Let \(P:\NN\to\mathsf{Prop}\). If
\[
  P(0)
  \quad\text{and}\quad
  \forall n: \NN,\, P(n)\implies P(S(n)),
\]
then \(P(n)\) holds for every \(n:\NN\).
\end{thm}

\begin{proof}
Let \(A=\{n:\NN\mid P(n)\}\). The hypotheses say that \(A\)
contains \(0\) and is closed under \(S\). Because the constructors in
\xref{def:natural-number} exhaust \(\NN\), we have
\(A=\NN\).
\end{proof}

\begin{rem}[Number Systems as Types and Sets]{number-systems-types-and-sets}
Throughout these notes, a symbol denoting a number system may represent either
a type of numbers or the set of all numbers of that type. Thus \(n:\NN\) is
the type-theoretic form, whereas \(n\in\NN\) is the corresponding
set-theoretic form. The intended interpretation is determined by context.
\end{rem}

\begin{conv}[Positive Natural Numbers]{positive-natural-numbers}
We write
\[
  \NNOne = \NN \setminus \{0\}
\]
when zero must be excluded, such as for one-based indices or denominators.
\end{conv}

\begin{dfn}[Integers]{integer}
The \emph{integers} are the whole numbers, including the positive integers,
the negative integers, and zero:
\[
  \ZZ=\{\ldots,-2,-1,0,1,2,\ldots\}.
\]
\end{dfn}

% With no explicit ID, this automatically creates def:rational-numbers.
\begin{dfn}[Rational Numbers]
A real number is called \emph{rational} if it can be expressed as a quotient
of two integers. The set of rational numbers is
\[
  \QQ
  =\left\{\frac{p}{q}:p,q\in\ZZ,\ q\neq0\right\}.
\]
\end{dfn}

\begin{prop}[Unique Reduced Form]{unique-reduced-rational-form}
For every \(x\in\QQ\), there exists a unique pair
\(p,q\in\ZZ\) such that
\[
  x=\frac{p}{q},
  \qquad q>0,
  \qquad \gcd(\abs{p},q)=1.
\]
\end{prop}

\begin{rem}{number-set-inclusions}
Under the usual identifications,
\(\NN\subset\ZZ\subset\QQ\).
\end{rem}

\begin{dfn}[Floor]{floor}
  For \(x \in \RR\), the \emph{floor} \(\lfloor x \rfloor\) is the unique \(m \in \ZZ\) satisfying
  \[
    m \leq x < m + 1.
  \]
\end{dfn}

\begin{dfn}[Ceiling]{ceiling}
  For \(x \in \RR\), the \emph{ceiling} \(\lceil x \rceil\) is the unique \(m \in \ZZ\) satisfying
  \[
    m - 1 < x \leq m.
  \]
\end{dfn}

\begin{thm}[Well-Definedness of Floor and Ceiling]{floor-ceiling-well-defined}
  For every \(x \in \RR\):
  \begin{itemize}[nosep]
    \item there exists a unique \(m \in \ZZ\) satisfying \(m \leq x < m + 1\), namely \(\lfloor x \rfloor\);
    \item there exists a unique \(m \in \ZZ\) satisfying \(m - 1 < x \leq m\), namely \(\lceil x \rceil\).
  \end{itemize}
\end{thm}

\begin{proof}
\missingprereq{Existence and uniqueness are deferred. Sketch: existence follows from the \xref[Archimedean property]{lec02::thm:archimedean-property} combined with the well-ordering of \(\NN\) (equivalent to \xref[Natural Number Induction]{thm:natural-number-induction}); uniqueness is immediate from the strict inequalities. A rigorous proof belongs in abstract algebra / order-theory prerequisites.}
\end{proof}

\section{Fields}

\begin{dfn}[Closed Operation]
An operation $\ast$ is said to be \emph{closed} on a set $A$ if for all $a,b \in A$, we have $a \ast b \in A$. 
\end{dfn}

\begin{dfn}[Field]{field}
A \emph{field} is a set \(F\) equipped with addition and multiplication such that both operations are
\begin{itemize}[nosep]
  \item \xref[closed operations]{def:closed-operation} on \(F\),
  \item associative, meaning for all \(a,b,c \in F\), we have \((a + b) + c = a + (b + c)\) and \((a \cdot b) \cdot c = a \cdot (b \cdot c)\), 
  \item commutative, meaning for all \(a,b \in F\), we have \(a + b = b + a\) and \(a \cdot b = b \cdot a\), 
  \item have distinct identity elements \(0,1\in F\), meaning
    \(a+0=a\), \(a\cdot1=a\), and \(0\neq1\),
  \item each element has an additive inverse, such that for each \(a \in F\), there exists \(-a \in F\) with \(a + (-a) = 0\), 
  \item each non-zero element has a multiplicative inverse, such that for each \(a \in F\) with \(a \neq 0\), there exists \(a^{-1} \in F\) with \(a \cdot a^{-1} = 1\), and
  \item multiplication distributes over addition, meaning
    \(a\cdot(b+c)=a\cdot b+a\cdot c\).
\end{itemize}
\end{dfn}

\begin{prop}[\(\QQ\) Is a Field]{rationals-are-a-field}
The rational numbers with their usual addition and multiplication form a
\xref[field]{def:field}.
\end{prop}

\begin{rem}[\(\NN\) and \(\ZZ\) Are Not Fields]{natural-and-integers-not-fields}
The natural numbers and the integers are not fields because they do not have
multiplicative inverses for every non-zero element. 

For a counter example, consider the integer \(2\). There is no integer \(x\) such that \(2 \cdot x = 1\) 
\end{rem}

\section{Ordered Set} 

\begin{dfn}[Transitive]{transitive}
An ordering on a set is \emph{transitive} if 
\[ \forall x, y, z \in F, \text{ if } x \le y \text{ and } y \le z, \text{ then } x \le z \]
\end{dfn}

\begin{dfn}[Reflexive]{reflexive}
An ordering on a set is \emph{reflexive} if
\[ \forall x \in F, x \leq x \]
\end{dfn}

\begin{dfn}[Antisymmetric]{antisymmetric}
An ordering on a set is \emph{antisymmetric} if
\[ \forall x, y \in F, (x \le y) \land (y \le x) \implies x = y \]
\end{dfn} 

\begin{dfn}[Partial Order]{partial-order}
A binary relation \(\leq\) on a set \(F\) is a \emph{partial order} if it is
\xref[reflexive]{def:reflexive}, \xref[antisymmetric]{def:antisymmetric}, and
\xref[transitive]{def:transitive}.
\end{dfn}

\begin{dfn}[Totality]{totality}
An ordering on a set \(F\) is \emph{total} if every pair of elements is
comparable:
\[
  \forall x,y\in F,\quad x\leq y\lor y\leq x.
\]
\end{dfn}

\begin{dfn}[Total Order]{total-order}
  A \xref[partial order]{def:partial-order} \(\leq\) on \(F\) is a
  \emph{total order} if it also satisfies \xref[totality]{def:totality}.
\end{dfn}

\begin{conv}[Order Notation]{order-notation}
For \(a,b\in F\), define
\[
  a\geq b\iff b\leq a,\qquad
  a<b\iff a\leq b\land a\neq b,\qquad
  a>b\iff b<a.
\]
\end{conv}

\begin{dfn}[Binary Maximum]{binary-max}
  Let \(X\) be equipped with a \xref[total order]{def:total-order} \(\leq\). For \(a, b \in X\), define
  \[
    \max(a, b) :=
    \begin{cases}
      a & \text{if } a \geq b, \\
      b & \text{otherwise.}
    \end{cases}
  \]
\end{dfn}

\begin{dfn}[Binary Minimum]{binary-min}
  Let \(X\) be equipped with a \xref[total order]{def:total-order} \(\leq\). For \(a, b \in X\), define
  \[
    \min(a, b) :=
    \begin{cases}
      a & \text{if } a \leq b, \\
      b & \text{otherwise.}
    \end{cases}
  \]
\end{dfn}

\begin{dfn}[Maximum of a Set]{set-maximum}
  Let \(X\) be equipped with a \xref[total order]{def:total-order} \(\leq\), and let \(A \subseteq X\) be nonempty. An element \(m \in A\) is a \emph{maximum of \(A\)}, written \(m = \max A\), if
  \[
    \forall x \in A,\, m \geq x.
  \]
\end{dfn}

\begin{dfn}[Minimum of a Set]{set-minimum}
  Let \(X\) be equipped with a \xref[total order]{def:total-order} \(\leq\), and let \(A \subseteq X\) be nonempty. An element \(m \in A\) is a \emph{minimum of \(A\)}, written \(m = \min A\), if
  \[
    \forall x \in A,\, m \leq x.
  \]
\end{dfn}

\begin{prop}[Finite Nonempty Families Have a Maximum and Minimum]{finite-max-min-inductive}
  Let \(X\) be a totally ordered set. For any nonempty finite family \((a_1, \ldots, a_n)\) in \(X\), the maximum \(\max\{a_1, \ldots, a_n\}\) and minimum \(\min\{a_1, \ldots, a_n\}\) exist. They may be computed inductively:
  \[
    \max\{a_1\} = a_1, \qquad \max\{a_1, \ldots, a_k\} = \max\bigl(\max\{a_1, \ldots, a_{k-1}\},\, a_k\bigr) \text{ for } k \geq 2,
  \]
  and symmetrically for \(\min\).
\end{prop}

\begin{dfn}
[Ordered Field]{ordered-field}
  An \emph{ordered field} is a \xref[field]{def:field} \(F\) equipped with a
  \xref[total order]{def:total-order} \(\leq\) such that for all
  \(x,y,z\in F\),
  \begin{itemize}[nosep]
    \item if \(x\leq y\), then \(x+z\leq y+z\); and
    \item if \(0\leq x\) and \(0\leq y\), then \(0\leq xy\).
  \end{itemize} 
\end{dfn}

\begin{dfn}[Absolute Value]{absolute-value}
  Let \(F\) be an \xref[ordered field]{def:ordered-field}. For \(a\in F\),
  the \emph{absolute value} of \(a\) is
  \[
    \abs{a}:=
    \begin{cases}
      a,  & a\geq0,\\
      -a, & a<0.
    \end{cases}
  \]
\end{dfn}

\begin{lem}[Difference Characterization]{difference-characterization}
  Let \(F\) be an ordered field. For all \(a,b\in F\),
  \[
    a<b\iff 0<b-a.
  \]
\end{lem}

\begin{proof}
  Since addition preserves order and is cancellative in an ordered field,
  \[
    a<b
    \iff a-a<b-a
    \iff 0<b-a.
  \]
  Here \(a-a=0\) by the additive-inverse axiom in the
  \xref[definition of a field]{def:field}.
\end{proof}

\begin{lem}[Positive Perturbation]{positive-perturbation}
  Let \(F\) be an ordered field. For all \(s,\varepsilon\in F\),
  \[
    \varepsilon>0
    \implies
    s-\varepsilon<s<s+\varepsilon.
  \]
\end{lem}

\begin{proof}
  If \(\varepsilon>0\), then
  \[
    s-(s-\varepsilon)=\varepsilon>0
    \quad\text{and}\quad
    (s+\varepsilon)-s=\varepsilon>0.
  \]
  Both inequalities follow from
  \xref[the difference characterization]{lem:difference-characterization}.
\end{proof}

\section{Indexed Family}

\begin{dfn}[Indexed Family]{indexed-family}
  Let \(I\) be an index set or type, and let \(X\) be a set or type. An
  \emph{\(I\)-indexed family in \(X\)} is a mapping
  \[
    a:I\to X.
  \]
\end{dfn}

\begin{dfn}[Term]{term}
  Let \(a:I\to X\) be an indexed family. For each \(i\in I\), the value
  \[
    a_i:=a(i)
  \]
  is called the \emph{term indexed by \(i\)}. For a sequence, \(a_n\) is
  called the \emph{\(n\)-th term}.
\end{dfn}

\begin{conv}[Indexed Family Notation]{indexed-family-notation}
  We denote an indexed family \(a:I\to X\) by
  \[
    (a_i)_{i\in I}.
  \]
  We may write simply \((a_i)\) when the index set is clear.
\end{conv}

\begin{dfn}[Sequence]{sequence}
  Let \(I=\NN\) or \(I=\NNOne\). An \(I\)-indexed family
  \(a:I\to X\) is called a \emph{sequence in \(X\)}. For a sequence, we
  conventionally use \(n\) as the index and write \((a_n)_{n\in I}\).
\end{dfn}

\begin{eg}
  \(\left(\frac{1}{n}\right)_{n\in\NNOne}\) is a sequence.
\end{eg}

\section{Finite Sums and Products}

\begin{dfn}[Finite Summation]{finite-summation}
  Let \(F\) be a \xref[field]{def:field}, \(a : \NNOne \to F\), and \(m, n \in \NNOne\). The \emph{finite summation} \(\sum_{k = m}^{n} a_k\) is defined recursively:
  \[
    \sum_{k = m}^{n} a_k =
    \begin{cases}
      0 & \text{if } m > n, \\
      \left(\sum_{k = m}^{n - 1} a_k\right) + a_n & \text{if } m \leq n.
    \end{cases}
  \]
\end{dfn}

\begin{dfn}[Finite Product]{finite-product}
  Let \(F\) be a \xref[field]{def:field}, \(a : \NNOne \to F\), and \(m, n \in \NNOne\). The \emph{finite product} \(\prod_{k = m}^{n} a_k\) is defined recursively:
  \[
    \prod_{k = m}^{n} a_k =
    \begin{cases}
      1 & \text{if } m > n, \\
      \left(\prod_{k = m}^{n - 1} a_k\right) \cdot a_n & \text{if } m \leq n.
    \end{cases}
  \]
\end{dfn}

\begin{prop}[Linearity of Summation]{summation-linearity}
  Let \(F\) be a \xref[field]{def:field}, \(a, b : \NNOne \to F\), \(c, d \in F\), and \(m, n \in \NNOne\). Then
  \[
    \sum_{k = m}^{n} (c \cdot a_k + d \cdot b_k)
    = c \cdot \sum_{k = m}^{n} a_k + d \cdot \sum_{k = m}^{n} b_k.
  \]
\end{prop}

\begin{prop}[Index Shift for Summation]{summation-shift}
  Let \(F\) be a \xref[field]{def:field}, \(a : \NNOne \to F\), and \(m, n, p \in \NNOne\) with \(m + p, n + p \in \NNOne\). Then
  \[
    \sum_{k = m}^{n} a_k = \sum_{k = m + p}^{n + p} a_{k - p}.
  \]
\end{prop}

\begin{prop}[Splitting of Summation]{summation-split}
  Let \(F\) be a \xref[field]{def:field}, \(a : \NNOne \to F\), and \(m, p, n \in \NNOne\) with \(m \leq p \leq n\). Then
  \[
    \sum_{k = m}^{n} a_k = \sum_{k = m}^{p} a_k + \sum_{k = p + 1}^{n} a_k.
  \]
\end{prop}

\begin{prop}[Multiplicativity of Product]{product-multiplicativity}
  Let \(F\) be a \xref[field]{def:field}, \(a, b : \NNOne \to F\), and \(m, n \in \NNOne\). Then
  \[
    \prod_{k = m}^{n} (a_k \cdot b_k) = \left(\prod_{k = m}^{n} a_k\right) \cdot \left(\prod_{k = m}^{n} b_k\right).
  \]
\end{prop}

\begin{prop}[Index Shift for Product]{product-shift}
  Let \(F\) be a \xref[field]{def:field}, \(a : \NNOne \to F\), and \(m, n, p \in \NNOne\) with \(m + p, n + p \in \NNOne\). Then
  \[
    \prod_{k = m}^{n} a_k = \prod_{k = m + p}^{n + p} a_{k - p}.
  \]
\end{prop}

\begin{prop}[Splitting of Product]{product-split}
  Let \(F\) be a \xref[field]{def:field}, \(a : \NNOne \to F\), and \(m, p, n \in \NNOne\) with \(m \leq p \leq n\). Then
  \[
    \prod_{k = m}^{n} a_k = \left(\prod_{k = m}^{p} a_k\right) \cdot \left(\prod_{k = p + 1}^{n} a_k\right).
  \]
\end{prop}

\section{Sets}

\begin{rem}[Membership and Classical Logic]{membership-classical-logic}
  For every set \(A\) and every \(x\) in the chosen universe, exactly one of
  \(x\in A\) and \(x\notin A\) holds.
\end{rem}

\begin{proof}
  Let \(P\) be the proposition \(x\in A\).
  By the law of excluded middle, \(P\lor\neg P\), while the law of
  non-contradiction gives \(\neg(P\land\neg P)\).
  Hence exactly one of \(x\in A\) and \(x\notin A\) holds.
\end{proof}

\begin{dfn}[Intersection]{intersection}
  For sets \(A\) and \(B\),
  \[
    x\in A\cap B \iff x\in A\text{ and }x\in B.
  \]
\end{dfn}

\begin{dfn}[Indexed Intersection]{indexed-intersection}
  Let \(I\) be a nonempty index type and let
  \[
    A:I\to\mathbf{Set}
  \]
  be an \xref[indexed family]{def:indexed-family} of sets. The
  \emph{indexed intersection} is characterized by
  \[
    x\in\bigcap_{i:I}A_i
    \iff
    \forall i:I,\quad x\in A_i.
  \]
  In particular, for \(a,b\in\NN\) with \(a\leq b\),
  \[
    \bigcap_{n=a}^{b}A_n
    \equiv
    \bigcap_{n\in\{m\in\NN:a\leq m\leq b\}}A_n,
    \qquad
    \bigcap_{n=0}^{\infty}A_n
    \equiv
    \bigcap_{n:\NN}A_n.
  \]
\end{dfn}

\begin{dfn}[Union]{union}
  For sets \(A\) and \(B\),
  \[
    x\in A\cup B \iff x\in A\text{ or }x\in B.
  \]
\end{dfn}

\begin{dfn}[Indexed Union]{indexed-union}
  Let \(I\) be an index type and let
  \[
    A:I\to\mathbf{Set}
  \]
  be an \xref[indexed family]{def:indexed-family} of sets. The
  \emph{indexed union} is characterized by
  \[
    x\in\bigcup_{i:I}A_i
    \iff
    \exists i:I,\quad x\in A_i.
  \]
  In particular, for \(a,b\in\NN\) with \(a\leq b\),
  \[
    \bigcup_{n=a}^{b}A_n
    \equiv
    \bigcup_{n\in\{m\in\NN:a\leq m\leq b\}}A_n,
    \qquad
    \bigcup_{n=0}^{\infty}A_n
    \equiv
    \bigcup_{n:\NN}A_n.
  \]
  If \(I\) is empty, then \(\bigcup_{i:I}A_i=\emptyset\).
\end{dfn}

\begin{dfn}[Complement]{complement}
  Given a universal set \(U\) and \(A\subseteq U\), the complement of \(A\)
  is \(A^c=U\setminus A\); equivalently,
  \[
    x\in A^c \iff x\notin A.
  \]
\end{dfn}

\begin{dfn}[Set Difference]{set-difference}
  For sets \(A\) and \(B\),
  \[
    x\in A\setminus B \iff x\in A\text{ and }x\notin B.
  \]
\end{dfn}

\begin{dfn}[Subset]{subset}
  For sets \(A\) and \(B\),
  \[
    A\subseteq B \iff \text{for every }x,\quad x\in A\implies x\in B.
  \]
\end{dfn}

\begin{dfn}[Set Equality]{set-equality}
  For sets \(A\) and \(B\),
  \[
    A=B \iff A\subseteq B\text{ and }B\subseteq A.
  \]
\end{dfn} 

\begin{theorem}[De Morgan's Laws]{de-morgan-laws}
  For sets \(A\) and \(B\) in a universal set \(U\),
  \[
    (A\cap B)^c=A^c\cup B^c
    \quad\text{and}\quad
    (A\cup B)^c=A^c\cap B^c.
  \]
  Also, 
  \[
    \left(\bigcap_{n=1}^{\infty} A_n\right)^c = \bigcup_{n=1}^{\infty} (A_n)^c \quad \text{and} \quad \left(\bigcup_{n=1}^{\infty} A_n\right)^c = \bigcap_{n=1}^{\infty} (A_n)^c 
  \]
\end{theorem}

\begin{eg}[Nested Sequence]{nested-sequence} 
Let $A_n = \{n, n+1, n+2, \dots\}, n \in \NNOne$ 

ie, $A_1 = \{1, 2, 3, \dots\}, A_5 = \{5, 6, 7, \dots\}$ 

Trivially we can get $A_1 \subset A_2 \subset \dots $. This is called a nested sequence. 

Where 
\[
\bigcup^{\infty}_{n=1} A_n = A_1 \quad \bigcap^{\infty}_{n=1} = \emptyset 
\] 

The intersection is $\emptyset$ due to if a number, let's say for any natural number $p$ is in the set, then $A_{p+1}$ does not include that number. 

So in the universe $\NN$, 
\[\forall p \notin \bigcap^{\infty}_{n=1} \implies \bigcap^{\infty}_{n=1} = \emptyset \] 
\end{eg}

\section{Mapping} 

\begin{dfn}[Mapping]
  A mapping \(f: A \to B\) is a set of rule st.: 
  \[\forall x \in A, \exists! y = f(x) \in B \leftrightarrow \forall x \in A, \exists! y \in B \text{ st. } (x,y) \in f\]
\end{dfn} 

\begin{dfn}[Domain]
  The \emph{domain} of a mapping \(f: A \to B\) is the set \(A\).
\end{dfn} 

\begin{dfn}[Codomain]
  The \emph{codomain} of a mapping \(f: A \to B\) is the set \(B\).
\end{dfn} 

\begin{dfn}[Range] 
  The \emph{range} of a mapping \(f: A \to B\) is a set \(C\) such that: 
  \[
    C = \{y: \exists x \in A, \text{ st. } y = f(x)\} 
  \] 
\end{dfn} 

\begin{rem}[Range and codomains are different]
  The \xref[range]{def:range} is a subset of \xref[codomain]{def:codomain}. Notice that some of the values in the \xref[codomain]{def:codomain} is not reachable. 

  For example, 
  \[f = x^2: \RR \to \RR\] 

  The codomain is \(\RR\) but the range is \(\RR^+ \cup \{0\}\). 
\end{rem}

Let \(\mathbb B:=\{\mathrm{tt},\mathrm{ff}\}\) denote the set of Boolean values.

\begin{dfn}[Injective Mapping/One to One Mapping]{injective-mapping}
  A mapping \(f: A \to B\) is called \emph{injective} or \emph{one to one} if 
  \[
    f(x_1) = f(x_2) \implies x_1 = x_2
  \] 
  equivalently,
  \[
    x_1 \neq x_2 \implies f(x_1) \neq f(x_2).
  \]
\end{dfn} 

\begin{eg}
  The following is an injective mapping:
  \[
    f:\mathbb B\to\ZZ,\qquad
    \mathrm{tt}\mapsto 0,\quad \mathrm{ff}\mapsto 1.
  \]
\end{eg}
  
\begin{dfn}[Surjective Mapping/onto Mapping]{surjective-mapping}
  A mapping \(f: A \to B\) is called \emph{surjective} or \emph{onto} if 
  \[
    \forall y \in B, \exists x \in A, \text{ st. } f(x) = y
  \] 
  equivalently, 
  \[
    \text{Range}(f) = B 
  \]
\end{dfn} 

\begin{eg}
  The following is a surjective mapping:
  \[
    g:\ZZ\to\mathbb B,\qquad
    g(n)=
    \begin{cases}
      \mathrm{tt}, & n=0,\\
      \mathrm{ff}, & n\neq 0.
    \end{cases}
  \]
\end{eg}

\begin{dfn}[Bijective Mapping]{bijective-mapping}
  A mapping \(f: A \to B\) is called \emph{bijective} if it is both
  \xref[injective]{def:injective-mapping} and
  \xref[surjective]{def:surjective-mapping}. 
\end{dfn} 

\begin{eg}
  The following is a bijective mapping, so it is also injective and surjective:
  \[
    h:\{0,1\}\to\mathbb B,\qquad
    0\mapsto\mathrm{tt},\quad 1\mapsto\mathrm{ff}.
  \]
\end{eg}

\section{Distance}

\begin{dfn}[Distance]{distance}
  Let \(F\) be an \xref[ordered field]{def:ordered-field}. For \(x, y \in F\), the \emph{distance} between \(x\) and \(y\) is
  \[
    d(x, y) := \abs{x - y}.
  \]
\end{dfn}

\begin{lem}[Symmetry of Distance]{distance-symmetry}
  Let \(F\) be an \xref[ordered field]{def:ordered-field}. For all \(x, y \in F\),
  \[
    \dist{x}{y} = \dist{y}{x}.
  \]
\end{lem}

\begin{proof}
  By \xref[definition]{def:distance}, \(\dist{x}{y} = \abs{x - y}\) and \(\dist{y}{x} = \abs{y - x} = \abs{-(x - y)}\). It suffices to show \(\abs{-(x - y)} = \abs{x - y}\); we split on the sign of \(x - y\).

  Case 1: \(x - y \geq 0\). Then \(\abs{x - y} = x - y\). Also \(-(x - y) \leq 0\), so \(\abs{-(x - y)} = -(-(x - y)) = x - y = \abs{x - y}\).

  Case 2: \(x - y < 0\). Then \(\abs{x - y} = -(x - y)\). Also \(-(x - y) > 0\), so \(\abs{-(x - y)} = -(x - y) = \abs{x - y}\).
\end{proof}

\begin{lem}[Triangle Inequality for Distance]{triangle-inequality-distance}
  Let \(F\) be an \xref[ordered field]{def:ordered-field}. For all \(x, y, z \in F\),
  \[
    \dist{x}{z} \leq \dist{x}{y} + \dist{y}{z}.
  \]
\end{lem}

\begin{proof}
  Let \(a = x - y\) and \(b = y - z\). Then \(x - z = a + b\), so \(\dist{x}{z} = \abs{a + b}\).
  We split on the sign of \(a + b\).

  Case 1: \(a + b \geq 0\). Then \(\abs{a + b} = a + b\). Since \(a \leq \abs{a}\) and \(b \leq \abs{b}\),
  \[
    \abs{a + b} = a + b \leq \abs{a} + \abs{b} = \dist{x}{y} + \dist{y}{z}.
  \]

  Case 2: \(a + b < 0\). Then \(\abs{a + b} = -(a + b) = -a - b\). Since \(-a \leq \abs{a}\) and \(-b \leq \abs{b}\),
  \[
    \abs{a + b} = -a - b \leq \abs{a} + \abs{b} = \dist{x}{y} + \dist{y}{z}.
  \]
\end{proof}

\begin{lem}[\(\varepsilon\)-Criterion for Equality]{epsilon-equality}
  Let \(F\) be an \xref[ordered field]{def:ordered-field}. For \(a, b \in F\),
  \[
    a=b
    \iff
    \forall \varepsilon \in F,\, \varepsilon > 0,\, \abs{a-b}<\varepsilon.
  \]
\end{lem}

\begin{proof}In order to proof \xref[lem:epsilon-equality]{lem:epsilon-equality}, we need to prove both directions. 

\noindent\textbf{(a)}: $a=b \implies \forall \varepsilon>0, \abs{a-b}<\varepsilon$

Suppose \(a=b\). Then \(\abs{a-b}=0\). For any \(\varepsilon>0\), we have \(\abs{a-b}=0<\varepsilon\).

\noindent\textbf{(b)}: $\forall \varepsilon>0, \abs{a-b}<\varepsilon \implies a=b$

Assume \(a \neq b\). 

Then, this means that \[\exists \varepsilon_0 > 0, \abs{a - b} = \varepsilon_0 \]

Taking \(\varepsilon=\varepsilon_0\) gives \(\varepsilon_0=\abs{a-b}<\varepsilon_0\), a contradiction.

So, we have \(\forall \varepsilon>0, \abs{a-b}<\varepsilon \implies \lnot a \neq b\)

Since $\lnot a \neq b \iff a = b$, So, $\forall \varepsilon>0, \abs{a-b}<\varepsilon \implies a=b$

Therefore, both implications hold, and hence proven. 
\end{proof} 

\section{Bounds} 

\begin{dfn}[Bounded Above]
  A set \(A\) is \emph{bounded above} in the universe $U$ if \[\exists b \in U, \forall a \in A, a \leq b\]  
  Such $b$ is called an \emph{upper bound} of \(A\). 
\end{dfn} 

\begin{dfn}[Bounded Below]
  A set \(A\) is \emph{bounded below} in the universe $U$ if \[\exists b \in U, \forall a \in A, a \geq b\] 
  Such $b$ is called a \emph{lower bound} of \(A\). 
\end{dfn} 

\begin{dfn}[Least Upper Bound/Supremum]{supremum} 
  $s$ is the \emph{least upper bound} or \emph{supremum} of a set \(A\) if 
  \begin{enumerate}[nosep]
    \item $s$ is an upper bound of \(A\)
    \item $\forall b$ is an upper bound of \(A\), $s \leq b$ 
  \end{enumerate}
\end{dfn} 

\begin{prop}[\(\varepsilon\)-Criterion for the Supremum]
{epsilon-supremum}
  Let \(F\) be an ordered field, let \(A\subseteq F\) be nonempty, and
  let \(s\in F\).
  Then \(s=\sup A\) if and only if
  \begin{enumerate}[nosep]
    \item \(s\) is an upper bound of \(A\); and
    \item for every \(\varepsilon\in F\) with \(\varepsilon>0\),
      \(s-\varepsilon\) is not an upper bound of \(A\).
  \end{enumerate}
  Equivalently,
  \[
    s=\sup A
    \iff
    \begin{cases}
      \forall x\in A,\quad x\leq s,\\
      \forall\varepsilon\in F,\quad \varepsilon>0\implies
      \exists x\in A\text{ s.t. }s-\varepsilon<x.
    \end{cases}
  \]
\end{prop} 

\begin{proof}
In order to proof \(\varepsilon\)-Characterization of the Supremum, we need to proof both directions are true. 

\textbf{(a)}: 
  \[
    s=\sup A
    \implies
    \begin{cases}
      \forall x\in A,\quad x\leq s,\\
      \forall\varepsilon>0,\quad
      \exists x\in A\text{ s.t. }s-\varepsilon<x.
    \end{cases}
  \] 

  According to the \xref[definition of supremum]{def:supremum}, we know that $s$ is an upper bound of \(A\), ie, $\forall x\in A, x \leq s$. 

  Take arbitrary \(\varepsilon\in F\) with \(\varepsilon>0\). By \xref{lem:positive-perturbation},
  \(
  s' := s-\varepsilon, \, s'<s.
  \)

  Again, according to the \xref[definition of supremum]{def:supremum}, we know $\forall b \text{ is an upper bound of } A, s \leq b$.
  Since $s' < s$, we have that $s'$ is not an upper bound of $A$. 

  Therefore, $s' = s - \varepsilon$ is not an upper bound of $A$. According to the \xref[definition of upper bound]{def:bounded-above}, $\exists x\in A, \, x > s' \implies \exists x > s - \varepsilon$. 

  So, (a) part is proven. 

\textbf{(b)}: 
  \[
    \begin{cases}
      \forall x\in A,\quad x\leq s,\\
      \forall\varepsilon>0,\quad
      \exists x\in A\text{ s.t. }s-\varepsilon<x.
    \end{cases}
    \implies
    s=\sup A
  \] 

  Assume $\forall x\in A,\quad x\leq s$ and \(\forall\varepsilon>0,\quad \exists x\in A\text{ s.t. }s-\varepsilon<x\), and suppose for contradiction that $s \neq \sup A$. 

  From $\forall x\in A,\quad x\leq s$, it follows that $s$ is an upper bound of \(A\). 

  If $s \neq \sup A$, then there exists an upper bound $b$ of \(A\) such that $b < s$. Let
  \(
    \varepsilon := s-b > 0.
  \)
  By the assumption, there exists $x\in A$ such that
  \(
    s-\varepsilon<x.
  \)
  However,
  \(
    s-\varepsilon=s-(s-b)=b,
  \) 
  so $b<x$, contradicting that $b$ is an upper bound of $A$.

  So (b) is proven. 
\end{proof}

\begin{dfn}[Greatest Lower Bound/Infimum]{infimum} 
  $i$ is the \emph{greatest lower bound} or \emph{infimum} of a set \(A\) if 
  \begin{enumerate}[nosep]
    \item $i$ is a lower bound of \(A\)
    \item $\forall b$ is a lower bound of \(A\), $i \geq b$ 
  \end{enumerate}
\end{dfn} 

\begin{prop}[\(\varepsilon\)-Criterion for the Infimum]
{epsilon-infimum}
  Let \(F\) be an ordered field, let \(A\subseteq F\) be nonempty, and
  let \(i\in F\).
  Then \(i=\inf A\) if and only if
  \begin{enumerate}[nosep]
    \item \(i\) is a lower bound of \(A\); and
    \item for every \(\varepsilon\in F\) with \(\varepsilon>0\),
      \(i+\varepsilon\) is not a lower bound of \(A\).
  \end{enumerate}
  Equivalently,
  \[
    i=\inf A
    \iff
    \begin{cases}
      \forall x\in A,\quad i\leq x,\\
      \forall\varepsilon\in F,\quad \varepsilon>0\implies
      \exists x\in A\text{ s.t. }x<i+\varepsilon.
    \end{cases}
  \]
\end{prop}

\begin{conv}
  We use the notation $\sup A$ to denote the supremum of a set \(A\) and $\inf A$ to denote the infimum of a set \(A\).
\end{conv} 

\begin{rem}[Supremums and Infimums are unique]{supremum-infimum-uniqueness}
  If a set \(A\) has a supremum, then it is unique. If a set \(A\) has an infimum, then it is unique. 
\end{rem} 

\begin{proof}[Supremums are unique]
  Assume there are 2 supremums for set \(A\), \(s_1, s_2\). Where $s_1 \neq s_2$. 

  Since $\forall s$ is an upper bound of \(A\), $s_1 \leq s$, also $s_2 \leq s$. 

  Notice that $s_1$ is an upper bound of \(A\), and $s_2$ is a least upper bound, so $s_2 \leq s_1$. 

  Also, $s_2$ is an upper bound of \(A\), and $s_1$ is a least upper bound, so $s_1 \leq s_2$. 

  Trivially, $s_1 = s_2$. Contradicts with $s_1 \neq s_2$. 

  So there are no multiple supremum. Supremums are unique for each set. 
\end{proof} 

\begin{dfn}[Bounded Set]
  A set \(A\) is \emph{bounded} in the universe $U$ if it is both \xref[bounded above]{def:bounded-above} and \xref[bounded below]{def:bounded-below}. 
\end{dfn} 


\section{Axiom of Completeness} 

\begin{dfn}[Complete Ordered Field]{complete-ordered-field}
  A \emph{complete ordered field} is an ordered field \(F\) such that every
  nonempty set \(A\subseteq F\) that is
  \xref[bounded above]{def:bounded-above} has a
  \xref[least upper bound]{def:supremum} in \(F\).
\end{dfn}

\begin{rem}[Hierarchy of Order Structures]{order-structure-hierarchy}
Just now discussed the four ordered structures, each strictly extends the previous one and unlocks one new capability. When stating a theorem, use the weakest tier that suffices.
\begin{itemize}[nosep]
  \item \xref[Partial order]{def:partial-order}: comparison via \(\leq\) (reflexive, antisymmetric, transitive). Minimal and maximal elements make sense when they exist.
  \item \xref[Total order]{def:total-order}: adds totality, so any two elements are comparable. This is the minimum needed to define \xref[intervals]{lec02::def:interval}.
  \item \xref[Ordered field]{def:ordered-field}: adds the field operations \(+,-,\cdot\) compatible with \(\leq\). Enables \xref[absolute value]{def:absolute-value}, distance, and \(\varepsilon\)-arguments.
  \item \xref[Complete ordered field]{def:complete-ordered-field}: adds the \xref[Axiom of Completeness]{axm:aoc}. Suprema and infima of bounded sets exist, and limit-existence results (e.g.\ NIP) require this tier.
\end{itemize}
\end{rem}

\begin{axm}[Axiom of Completeness]{aoc}
Every nonempty set \(A\subseteq\RR\) that is \xref[bounded above]{def:bounded-above} has a 
\xref[least upper bound]{def:supremum} in \(\RR\) 
\end{axm}

\begin{lem}[Supremum--Infimum Equivalence]{supremum-infimum-equivalence}
  Let \(F\) be an ordered field. The following statements are equivalent:
  \begin{enumerate}[nosep]
    \item Every nonempty set \(A\subseteq F\) that is bounded above in
      \(F\) has a supremum in \(F\).
    \item Every nonempty set \(A\subseteq F\) that is bounded below in
      \(F\) has an infimum in \(F\).
  \end{enumerate}
  Consequently, the \xref[definition of a complete ordered field]
  {def:complete-ordered-field} may equivalently use the infimum property.
  In particular, \xref[the Axiom of Completeness]{axm:aoc} is equivalent
  to the assertion that every nonempty subset of \(\RR\) bounded
  below has an infimum in \(\RR\).
\end{lem}

\begin{proof}
  First assume the supremum property, and let \(A\subseteq F\) be nonempty
  and bounded below. Define
  \[
    -A=\{-a:a\in A\}.
  \]
  Then \(-A\) is nonempty and bounded above, so let \(s=\sup(-A)\).
  For every \(a\in A\), we have \(-a\le s\), and hence \(-s\le a\).
  Thus \(-s\) is a lower bound of \(A\). If \(\ell\) is any lower bound
  of \(A\), then \(-\ell\) is an upper bound of \(-A\). Therefore
  \(s\le-\ell\), so \(\ell\le-s\). Hence \(-s=\inf A\).

  Conversely, assume the infimum property, and let \(A\subseteq F\) be
  nonempty and bounded above. Then \(-A\) is nonempty and bounded below,
  so let \(i=\inf(-A)\). For every \(a\in A\), we have \(i\le-a\), and
  hence \(a\le-i\). Thus \(-i\) is an upper bound of \(A\). If \(u\) is
  any upper bound of \(A\), then \(-u\) is a lower bound of \(-A\).
  Therefore \(-u\le i\), so \(-i\le u\). Hence \(-i=\sup A\).
\end{proof}

\begin{rem}[Integers and Natural Numbers]{integers-naturals-not-fields}
  Neither \(\ZZ\) nor \(\NN\) is a field, so neither can be
  called a complete ordered field. As ordered sets, however, their nonempty
  subsets that are bounded above have greatest elements.
\end{rem}

\begin{eg}
If $A = {\frac{1}{n}: n \in \NNOne}$, then $\sup A = 1$ and $\inf A = 0$. Although it has a infimum, it does not have a minimum since $0 \notin A$. 
\end{eg} 

\begin{cor}[Supremum and Infimum Existence in $\RR$]{supremum-infimum-existence}
  Every nonempty set \(A\subseteq\RR\) that is \xref[bounded above]{def:bounded-above} has a 
\xref[least upper bound]{def:supremum} in \(\RR\) and every nonempty set \(A\subseteq\RR\) that is \xref[bounded below]{def:bounded-below} has a 
\xref[greatest lower bound]{def:infimum} in \(\RR\) 
\end{cor}

\begin{proof}
  The supremum statement is \xref[the Axiom of Completeness]{axm:aoc},
  and the infimum statement follows from
  \xref{lem:supremum-infimum-equivalence}.
\end{proof}

\begin{rem}[Intuition for Completeness]
  We can have a set, let's say, \[S = \{r \in \QQ: r^2 < 3\}, \text{where } r \subset \QQ \subset \RR \]

  Although the set $S$ is bounded above in $\QQ$, it does not have a supremum in $\QQ$. According to the \xref[AoC]{axm:aoc}, we know that $\exists s \in \RR, s = \sup S$. This means we can find a number that can be described as the least upper bound a subset of rational numbers which is not in the rational numbers. 

\end{rem}

\section{Irrational Numbers}

\begin{lem}[Parity of a Square]{square-parity}
For every \(n\in\ZZ\),
\[
  n^2\text{ is even}\iff n\text{ is even},
\]
and equivalently,
\[
  n^2\text{ is odd}\iff n\text{ is odd}.
\]
\end{lem}

\begin{prop}[No Rational Square Root of \(2\)]{no-rational-square-root-of-two}
There is no \(q\in\QQ\) such that \(q^2=2\).
\end{prop}

\begin{proof}
Suppose, for contradiction, that some \(q\in\QQ\) satisfies
\(q^2=2\).

According to the \xref{prop:unique-reduced-rational-form}, and $q\in\QQ$ 
\[\exists p, r \in \ZZ, q = \frac{p}{r}, r > 0, \gcd(\abs{p}, r) = 1\]

Then we have
\[\left(\frac{p}{r}\right)^2 = 2 \implies p^2 = 2r^2\]

Since \(p^2\) is even, the \xref[square-parity lemma]{lem:square-parity}
implies that \(p\) is even. Therefore,
\[\exists k \in \ZZ,\quad p = 2k.\]

So, 
\[(2k)^2 = 2r^2 \implies 2k^2 = r^2\]

By the same argument, \(r\) is also even. Meaning that both \(p\) and \(r\) are even, which contradicts the assumption that \(\gcd(\abs{p}, r) = 1\).

So, it contradicts the assumption that \(q\in\QQ\) satisfies \(q^2=2\). Therefore, there is no \(q\in\QQ\) such that \(q^2=2\). 
\end{proof}

\begin{rem}[Larger Number System]{sqrt-two-larger-number-system}
Within \(\QQ\), the proposition only says that the equation \(q^2=2\)
has no solution \(q\in\QQ\). Writing
\(\sqrt{2}\notin\QQ\) presupposes a larger number system in which
\(\sqrt{2}\) exists. 
\end{rem}

\begin{rem}[Completeness]{rational-incompleteness}
The existence, in a larger number system, of some
\(x\notin\QQ\) satisfying \(x^2=2\) does not by itself prove that
\(\QQ\) is incomplete. Incompleteness must be demonstrated
internally by finding a nonempty subset of \(\QQ\) that is bounded above
in \(\QQ\) but has no supremum in \(\QQ\), as in
\xref[Proposition~\ref*{prop:rationals-not-complete}]
{prop:rationals-not-complete}.
\end{rem} 

\begin{conv}[Irrational Numbers]{irrational-numbers}
  We use \(\II\) to denote the set of \emph{irrational numbers} 
\end{conv} 

\begin{rem}[Number Sets]{number-sets-inclusions-real}
  The set of real numbers \(\RR\) is the union of the set of rational numbers \(\QQ\) and the set of irrational numbers \(\II\):
  \[
    \RR = \QQ \cup \II
  \]
  Also, in normal conventions, we have:
  \[
    \NN \subset \ZZ \subset \QQ \subset \RR
  \]
\end{rem} 

\begin{prop}[Integer--Irrational Root Dichotomy]{nth-root-dichotomy}
  For \(x,n\in\NN\) with \(n>1\), \(\sqrt[n]{x}\in\ZZ\) if
  \(x\) is a perfect \(n\)-th power, and otherwise
  \(\sqrt[n]{x}\in\II\).
\end{prop}

\begin{rem}[Sources of Irrational]
  \begin{itemize}[nosep]
    \item By \xref{prop:nth-root-dichotomy}, \(\sqrt[n]{x}\) is irrational
    whenever \(x\in\NN\), \(n>1\), and \(x\) is not a perfect
    \(n\)-th power.
    \item Taking logarithms 
    \item Constants defined by taking limits, such as $\pi$ and $e$. 
  \end{itemize} 
\end{rem} 

\begin{rem}[Irrational Numbers Are Not Closed]{irrational-not-closed}
  The irrational numbers are not closed under addition, multiplication, or
  exponentiation. For example,
  \[
    \sqrt{2}+(-\sqrt{2})=0,\qquad
    \sqrt{2}\cdot\sqrt{2}=2,\qquad
    (\sqrt{3})^{\log_3 4}=2,
  \]
  where all the inputs are irrational. Although \(\pi,e\in\II\), it
  is not known whether either \(\pi+e\) or \(\pi e\) is irrational; both are
  expected to be transcendental. By contrast, \(e^\pi\) has been proved
  transcendental, and is therefore irrational.
\end{rem}

\begin{prop}[The Rational Numbers Are Not Complete]
{rationals-not-complete}
  \(\QQ\) is not a
  \xref[complete ordered field]{def:complete-ordered-field}. In fact, the set
  \[
    A=\{q\in\QQ:q^2<2\}
  \]
  is nonempty and bounded above in \(\QQ\), but it has no supremum in \(\QQ\).
\end{prop}

\begin{proof}
  The set \(A\) is nonempty because \(1\in A\). It is bounded above in
  \(\QQ\) by \(2\): if some \(q\in A\) satisfied \(q>2\), then
  \(q^2>4>2\), contradicting \(q^2<2\).

  Suppose, for contradiction, that \(A\) has a supremum \(s\in\QQ\).
  Since \(1\in A\), we have \(s\geq1\).

  First suppose \(s^2<2\), and define
  \[
    h=\frac{2-s^2}{2s+2}\in\QQ.
  \]
  Then \(h>0\); moreover, \(s\geq1\) implies \(h<2\). Hence
  \begin{align*}
    (s+h)^2
      &=s^2+h(2s+h)\\
      &<s^2+h(2s+2)\\
      &=2.
  \end{align*}
  Thus \(s+h\in A\), but \(s+h>s\), contradicting that \(s\) is an
  upper bound of \(A\). Therefore \(s^2\not<2\).

  Now suppose \(s^2>2\), and define
  \[
    h=\frac{s^2-2}{2s}\in\QQ,
    \qquad
    t=s-h.
  \]
  Then \(h>0\), so \(t<s\), and
  \[
    t=\frac{s^2+2}{2s}>0,
    \qquad
    t^2=(s-h)^2=2+h^2>2.
  \]
  If some \(q\in A\) satisfied \(q>t\), then \(q>t>0\) and
  \[
    q^2-t^2=(q-t)(q+t)>0,
  \]
  so \(q^2>t^2>2\), contradicting \(q\in A\). Consequently every
  \(q\in A\) satisfies \(q\leq t\), making \(t\) an upper bound of
  \(A\). This contradicts the minimality of \(s\), because \(t<s\).
  Therefore \(s^2\not>2\).

  By totality, \(s^2=2\). This contradicts
  \xref[the fact that no rational number has square \(2\)]
  {prop:no-rational-square-root-of-two}. Hence \(A\) has no supremum in
  \(\QQ\), so \(\QQ\) is not complete.
\end{proof} 

\end{document}
  
